%==============================================================================
\section{Data Understanding --- Hiểu dữ liệu trước khi huấn luyện}
%==============================================================================

\begin{ynghia}[title={Nguồn}]
\tsurl{https://www.tutorialspoint.com/machine_learning/machine_learning_data_understanding.htm}.
\end{ynghia}

\begin{dongco}[title={Hai phần thường bị bỏ qua: toán và dữ liệu}]
Trong dự án ML, người ta thường bỏ qua hai phần quan trọng: \textbf{toán học} và \textbf{dữ liệu}.
Data understanding quan trọng vì ML là hướng \textbf{data-driven}: mô hình chỉ tốt bằng dữ liệu cung cấp.
\end{dongco}

\begin{dinhnghia}[title={Data understanding}]
Phân tích và khám phá dữ liệu để nhận diện pattern hoặc xu hướng có trong dữ liệu.
\end{dinhnghia}

\subsection{Các bước trong giai đoạn hiểu dữ liệu}

\begin{phuongphap}[title={Theo nguồn}]
\begin{itemize}
  \item \textbf{Data Collection}: thu thập dữ liệu liên quan từ database, website, API, \ldots
  \item \textbf{Data Cleaning}: loại dữ liệu không liên quan/trùng, xử lý missing; định dạng dễ phân tích.
  \item \textbf{Data Exploration}: khám phá pattern (histogram, scatter plot, correlation, \ldots).
  \item \textbf{Data Visualization}: biểu diễn trực quan (đồ thị, biểu đồ, bản đồ).
  \item \textbf{Data Preprocessing}: biến đổi dữ liệu cho phù hợp thuật toán ML (scale, đổi định dạng, giảm chiều).
\end{itemize}
\end{phuongphap}

\subsection{Vì sao phải hiểu dữ liệu trước khi đưa vào dự án ML?}

\begin{ynghia}[title={Bốn lý do chính}]
\begin{itemize}
  \item \textbf{Phát hiện vấn đề chất lượng}: missing, outlier, sai kiểu dữ liệu, không nhất quán --- xử lý sớm cải thiện mô hình.
  \item \textbf{Đánh giá mức độ liên quan}: biết feature nào quan trọng, feature nào có thể bỏ.
  \item \textbf{Chọn kỹ thuật ML phù hợp}: dữ liệu categorical $\Rightarrow$ classification; continuous $\Rightarrow$ regression, \ldots
  \item \textbf{Cải thiện hiệu năng mô hình}: feature engineering, tiền xử lý, chọn thuật toán đúng $\Rightarrow$ accuracy, precision, recall, F1 tốt hơn.
\end{itemize}
\end{ynghia}

\subsection{Hiểu dữ liệu bằng thống kê}

\begin{ynghia}[title={Hai hướng: thống kê và trực quan hoá}]
Sau bài nạp CSV, cần hiểu dữ liệu trước khi đưa vào mô hình.
Có thể hiểu bằng \textbf{thống kê} (bài này) hoặc trực quan hoá (Tutorialspoint có bài riêng ở phần sau).
Dưới đây là các recipe Python; dataset dùng \textbf{sklearn} (chạy được trên Linux/macOS/Windows).
\end{ynghia}

\begin{luuy}[title={Nguồn dữ liệu thay cho đường dẫn \texttt{C:\textbackslash...} trong bản gốc}]
Bản gốc Tutorialspoint dùng file CSV trên Windows. Ở đây dùng \texttt{fetch\_openml("diabetes")} (Pima) và \texttt{load\_iris()} (Iris) từ scikit-learn.
\end{luuy}

\subsubsection{Xem dữ liệu thô (Looking at Raw Data)}

\begin{vidu}[title={head() trên dataset Pima Indians Diabetes}]
\begin{lstlisting}
from sklearn.datasets import fetch_openml

names = ['preg', 'plas', 'pres', 'skin', 'test', 'mass', 'pedi', 'age', 'class']
data = fetch_openml("diabetes", version=1, as_frame=True, parser="auto").frame
data.columns = names
print(data.head(10))
\end{lstlisting}

\textbf{Output} (theo nguồn, 10 hàng đầu):
\begin{lstlisting}
preg   plas  pres    skin  test  mass   pedi    age      class
0      6      148     72     35   0     33.6    0.627    50    1
1      1       85     66     29   0     26.6    0.351    31    0
...
9      8      125     96      0   0     0.0     0.232    54    1
10     4      110     92      0   0     37.6    0.191    30    0
\end{lstlisting}
Cột đầu là chỉ số hàng, hữu ích để tham chiếu một quan sát cụ thể.
\end{vidu}

\subsubsection{Kích thước dữ liệu (Dimensions)}

\begin{ynghia}[title={Vì sao cần biết số hàng/cột}]
\begin{itemize}
  \item Quá nhiều hàng/cột $\Rightarrow$ huấn luyện lâu.
  \item Quá ít $\Rightarrow$ không đủ dữ liệu huấn luyện tốt.
\end{itemize}
\end{ynghia}

\begin{vidu}[title={shape trên Iris}]
\begin{lstlisting}
from sklearn.datasets import load_iris
import pandas as pd
data = pd.DataFrame(load_iris().data, columns=load_iris().feature_names)
print(data.shape)
\end{lstlisting}
\textbf{Output}: \texttt{(150, 4)} --- 150 hàng, 4 cột.
\end{vidu}

\subsubsection{Kiểu dữ liệu từng thuộc tính}

\begin{vidu}[title={dtypes trên Iris}]
\begin{lstlisting}
from sklearn.datasets import load_iris
import pandas as pd
data = pd.DataFrame(load_iris().data, columns=load_iris().feature_names)
print(data.dtypes)
\end{lstlisting}
\textbf{Output} (theo nguồn):
\begin{lstlisting}
sepal_length  float64
sepal_width   float64
petal_length  float64
petal_width   float64
dtype: object
\end{lstlisting}
Giúp biết khi nào cần chuyển kiểu (ví dụ string $\rightarrow$ float cho biến phân loại/thứ tự).
\end{vidu}

\subsubsection{Tóm tắt thống kê (describe)}

\begin{ynghia}[title={8 thống kê từ describe() cho mỗi cột số}]
Count, Mean, Standard Deviation, Min, Max, 25\%, Median (50\%), 75\%.
\end{ynghia}

\begin{vidu}[title={describe() trên Pima Indians Diabetes}]
\begin{lstlisting}
from sklearn.datasets import fetch_openml
from pandas import set_option

names = ['preg', 'plas', 'pres', 'skin', 'test', 'mass', 'pedi', 'age', 'class']
data = fetch_openml("diabetes", version=1, as_frame=True, parser="auto").frame
data.columns = names
set_option('display.width', 100)
set_option('precision', 2)
print(data.shape)
print(data.describe())
\end{lstlisting}
\textbf{Output}: shape \texttt{(768, 9)} và bảng thống kê mô tả đầy đủ các cột (count, mean, std, min, quartiles, max) theo nguồn.
\end{vidu}

\subsubsection{Phân phối lớp (Class Distribution)}

\begin{ynghia}[title={Quan trọng trong classification}]
Cần biết cân bằng lớp; lớp lệch mạnh có thể cần xử lý đặc biệt ở giai đoạn chuẩn bị dữ liệu.
\end{ynghia}

\begin{vidu}[title={groupby().size() trên cột class}]
\begin{lstlisting}
from sklearn.datasets import fetch_openml

names = ['preg', 'plas', 'pres', 'skin', 'test', 'mass', 'pedi', 'age', 'class']
data = fetch_openml("diabetes", version=1, as_frame=True, parser="auto").frame
data.columns = names
count_class = data.groupby('class').size()
print(count_class)
\end{lstlisting}
\textbf{Output} (theo nguồn):
\begin{lstlisting}
Class
0  500
1  268
dtype: int64
\end{lstlisting}
Lớp 0 gần gấp đôi lớp 1 --- dataset mất cân bằng lớp.
\end{vidu}

\subsubsection{Tương quan giữa các thuộc tính}

\begin{dinhnghia}[title={Correlation}]
Quan hệ giữa hai biến; hệ số tương quan Pearson phổ biến:
\begin{itemize}
  \item $= 1$: tương quan dương hoàn toàn.
  \item $= -1$: tương quan âm hoàn toàn.
  \item $= 0$: không tương quan tuyến tính.
\end{itemize}
\end{dinhnghia}

\begin{luuy}[title={Vì sao xem ma trận tương quan}]
Linear regression, logistic regression có thể kém nếu có cặp thuộc tính tương quan rất cao (đa cộng tuyến).
\end{luuy}

\begin{vidu}[title={corr() Pearson}]
\begin{lstlisting}
from sklearn.datasets import fetch_openml
from pandas import set_option

names = ['preg', 'plas', 'pres', 'skin', 'test', 'mass', 'pedi', 'age', 'class']
data = fetch_openml("diabetes", version=1, as_frame=True, parser="auto").frame
data.columns = names
set_option('display.width', 100)
set_option('precision', 2)
correlations = data.corr(method='pearson')
print(correlations)
\end{lstlisting}
Ma trận output cho tương quan giữa mọi cặp thuộc tính (theo bảng đầy đủ trong nguồn).
\end{vidu}

\subsubsection{Độ lệch phân phối (Skew)}

\begin{ynghia}[title={Skewness}]
Phân phối giả định Gaussian nhưng bị lệch trái/phải.
Nhiều thuật toán ML giả định dữ liệu gần phân phối chuẩn --- skew cao có thể cần hiệu chỉnh ở bước chuẩn bị dữ liệu.
\end{ynghia}

\begin{vidu}[title={skew() trên Pima dataset}]
\begin{lstlisting}
from sklearn.datasets import fetch_openml

names = ['preg', 'plas', 'pres', 'skin', 'test', 'mass', 'pedi', 'age', 'class']
data = fetch_openml("diabetes", version=1, as_frame=True, parser="auto").frame
data.columns = names
print(data.skew())
\end{lstlisting}
\textbf{Output} (theo nguồn):
\begin{lstlisting}
preg   0.90
plas   0.17
pres  -1.84
skin   0.11
test   2.27
mass  -0.43
pedi   1.92
age    1.13
class  0.64
dtype: float64
\end{lstlisting}
Giá trị gần 0 $\Rightarrow$ ít lệch; dương/âm lớn $\Rightarrow$ lệch rõ.
\end{vidu}
